\DocumentMetadata{
  pdfversion=2.0,
  pdfstandard=ua-2,
  lang=en-US,
  tagging=on,
  tagging-setup={math/setup=mathml-SE,tabsorder=structure}
}
\documentclass[11pt,letterpaper]{article}
\usepackage[margin=0.68in,includefoot]{geometry}
\usepackage{newcomputermodern}
\usepackage{amsmath,amscd,mathtools}
\providecommand{\square}{\mdlgwhtsquare}
\usepackage[hidelinks]{hyperref}
\hypersetup{
  pdftitle={Topology First-Year Exam, May 2019, Part 2},
  pdfauthor={University of Florida Department of Mathematics},
  pdfsubject={Graduate mathematics examination},
  pdfdisplaydoctitle=true,
  bookmarksopen=true
}
\setcounter{secnumdepth}{0}
\setlength{\parindent}{0pt}
\setlength{\parskip}{0.28em}
\setlength{\emergencystretch}{3em}
\setlength{\leftmargini}{2.15em}
\setlength{\leftmarginii}{2.65em}
\setlength{\leftmarginiii}{3em}
\renewcommand{\labelenumi}{\textbf{\theenumi.}}
\pagestyle{empty}
\raggedbottom
\allowdisplaybreaks
\newenvironment{examproblems}[1]{%
  \begin{enumerate}%
  \setcounter{enumi}{#1}%
  \setlength{\topsep}{0.35em}%
  \setlength{\partopsep}{0pt}%
  \setlength{\itemsep}{0.48em}%
  \setlength{\parsep}{0.18em}%
}{\end{enumerate}}
\newenvironment{examsubparts}{%
  \begin{enumerate}%
  \setlength{\topsep}{0.18em}%
  \setlength{\partopsep}{0pt}%
  \setlength{\itemsep}{0.16em}%
  \setlength{\parsep}{0.1em}%
}{\end{enumerate}}
\newenvironment{examinstructions}{%
  \par\begingroup\itshape
}{%
  \par\endgroup\vspace{0.18em}
}
\makeatletter
\renewcommand\section{\@startsection{section}{1}{0pt}%
  {-1.25ex plus -.25ex minus -.1ex}%
  {0.55ex plus .1ex}%
  {\normalfont\large\bfseries}}
\renewcommand{\@maketitle}{%
  \newpage\null\vskip -1.2em%
  {\footnotesize\bfseries
    \href{https://www.ufl.edu/}{University of Florida}, \href{https://math.ufl.edu/}{Department of Mathematics}\par}%
  \vskip 0.18em%
  \tagstructbegin{tag=Title}%
    {\LARGE\bfseries\href{https://gma.math.ufl.edu/exams/topology-first-year/}{Topology First-Year Exam}, May 2019, Part 2\par}%
  \tagstructend%
  \vskip 0.35em%
  \tagmcbegin{artifact}%
    \hrule height 1pt%
  \tagmcend%
  \vskip 0.55em%
}
\makeatother
\title{Topology First-Year Exam, May 2019, Part 2}
\author{}
\date{}
\begin{document}
\maketitle
\thispagestyle{empty}
\begin{examinstructions}
Answer all questions and work all problems. Make proofs as succinct as possible.
\end{examinstructions}
\begin{examproblems}{0}
\item
Let \(n \ge 1\). Suppose that \(f,g:X\to\mathbb{R}^n\) are continuous. Show that~\(f\) and~\(g\) are homotopic.
\item
Let \(\gamma:[0,1]\to X\) with basepoint \(x_0\). Let \(\gamma^{-1}\) be defined by \(\gamma^{-1}(t)=\gamma(1-t)\). Show that \(\gamma\ast\gamma^{-1}\) is loop homotopic to the constant loop \(1:[0,1]\to X\), \(1(t)\equiv x_0\).
\item
Show that \(\pi_1(S^1,1)=\mathbb{Z}\).
\item
Show that the disk, \(D^2=\{z\in\mathbb{C}\mid \lVert z\rVert\le 1\}\), has the fixed point property.
\item
Consider the matrix \(M=\begin{bmatrix}2&3\\1&1\end{bmatrix}\). This represents a group homomorphism \(M:\mathbb{Z}^2\to\mathbb{Z}^2\). Show that there is a map \(f:T^2\to T^2\) such that \(f_*=M\), where \(f_*:\pi_1(T^2)\to\pi_1(T^2)\) is the homomorphism induced by~\(f\) on the fundamental group.
\item
Suppose that \(f:S^1\to S^1\) is continuous. Show that~\(f\) is homotopic to \(z^n:S^1\to S^1\) for some \(n\in\mathbb{Z}\). This is the degree of~\(f\).
\item
Let \(n>1\). Let \(\mathbb{P}^n=S^n/D\), where~\(D\) identifies~\(x\) with its antipode for all \(x\in S^n\). Show that \(\pi_1(\mathbb{P}^n,x_0)\cong\mathbb{Z}_2\).
\item
Describe a space~\(X\) such that \(\pi_1(X,x_0)\cong\mathbb{Z}_3\). Verify that this is the fundamental group by the Seifert–van Kampen Theorem.
\item
A topological space is said to be Lindelöf provided that for every open cover~\(\mathcal{U}\) of~\(X\), there is a countable \(\mathcal{V}\subset\mathcal{U}\) covering~\(X\). Show that if~\(X\) is a separable metric space, then~\(X\) is Lindelöf.
\item
State and prove the Tychonoff Theorem.
\item
State the following theorems or verify the following statements with quick examples or statements.

\par
The Alexander Subbase Theorem.

\par
The Hahn–Mazurkiewicz Theorem

\par
The Brouwer Characterization of the Cantor Set

\par
The Jordan Curve Theorem

\par
The Borsuk–Ulam Theorem for \(S^2\).

\par
Determine the fundamental group of \(T^2\).

\par
Determine the fundamental group of the connected sum of two tori, \(\pi_1(T^2\#T^2)\).

\par
The Arcwise Connectedness Theorem

\par
Suppose that~\(X\) is a compact metric AR. Show that for any~\(Z\) and any pair of continuous functions \(f,g:Z\to X\), \(f\) and~\(g\) are homotopic.

\par
Suppose that~\(C\) is the Cantor set. Show that~\(C\) is homeomorphic to \(C\times C\).
\end{examproblems}
\end{document}
